\documentclass[11pt]{amsart}
\usepackage[T1]{fontenc}
\usepackage{amsmath,amssymb}
\usepackage[colorlinks=true,linkcolor=blue,citecolor=blue,urlcolor=blue]{hyperref}

\newcommand{\wstar}{w^*}

\begin{document}

\title{\(JL_2^*\) and Problem 1 of Barroso--Kalenda--Lin}
\author{Run fa\_banach\_001}
\date{2026-06-28}
\maketitle

\noindent
\textbf{Status.}
\texttt{literature\_implied\_answer (negative answer to Problem 1; Problem 2 not claimed)}.
This note records an existing literature implication, not a new theorem of the
run.

\section*{Source Question}

In \emph{On the approximate fixed point property in abstract spaces},
arXiv:1101.5274, Barroso--Kalenda--Lin ask the following as Problem~1:
for an arbitrary Hausdorff locally convex space \(X\), is it true that every
bounded sequence in \(X\) has a weakly Cauchy subsequence if and only if each
bounded separable subset of \(X\) is Frechet--Urysohn in the weak topology?
They point out that the answer is positive for locally convex spaces admitting
a compatible metrizable locally convex topology, and they suggest
\((X^*,w^*)\) for a Johnson--Lindenstrauss space as a possible source of a
negative answer.

\section*{Supporting Literature}

Martinez-Cervantes, \emph{Banach spaces with weak*-sequential dual ball},
arXiv:1612.05948, proves that the classical Johnson--Lindenstrauss space
\(JL_2\) has weak-star sequential dual ball of sequential order \(2\)
(Theorem~3.1 in that paper). The paper also recalls that weak-star
sequential dual ball implies weak-star sequential compactness of the dual
ball, and that \(JL_2\) does not have weak-star angelic dual. Equivalently,
\((B_{JL_2^*},w^*)\) is sequentially compact and sequential, but it is not
Frechet--Urysohn.

\section*{Identification}

Let \(Z=JL_2\), and set
\[
  E=(Z^*,\sigma(Z^*,Z)).
\]
This is a Hausdorff locally convex space. Its continuous dual is \(Z\), so the
weak topology of \(E\) is again \(\sigma(Z^*,Z)\).

First, every bounded sequence in \(E\) has a weakly Cauchy subsequence. Indeed,
if \((x_n^*)\) is bounded in \(E\), then it is pointwise bounded on \(Z\). By
the uniform boundedness principle, \((x_n^*)\) is norm bounded in the Banach
dual \(Z^*\). After rescaling, the sequence lies in \(B_{Z^*}\). Since
\((B_{Z^*},w^*)\) is weak-star sequentially compact, some subsequence converges
in \(\sigma(Z^*,Z)\). This subsequence is weakly convergent, hence weakly
Cauchy, in the locally convex space \(E\).

Second, not every bounded separable subset of \(E\) is Frechet--Urysohn in the
weak topology. Since \((B_{Z^*},w^*)\) is not Frechet--Urysohn, there are
\(A\subset B_{Z^*}\) and \(p\in \overline A^{w^*}\) such that no sequence in
\(A\) converges to \(p\). Since the dual ball is sequential, it has countable
tightness. Therefore there is a countable \(A_0\subset A\) with
\(p\in\overline{A_0}^{w^*}\). Let \(S=A_0\cup\{p\}\). The set \(S\) is
countable, hence separable, and is bounded in \(E\). In the subspace \(S\),
the point \(p\) lies in the closure of \(A_0\), but no sequence from \(A_0\)
converges to \(p\). Thus \(S\) is not Frechet--Urysohn.

Consequently the two properties in Problem~1 are not equivalent for arbitrary
Hausdorff locally convex spaces.

\section*{Scope}

This does not settle Problem~2 of Barroso--Kalenda--Lin, which asks whether
weak-AFPP is equivalent to the bounded-sequence weakly-Cauchy-subsequence
property for arbitrary Hausdorff locally convex spaces.

This packet is placed under literature-implied answers because arXiv:1612.05948
does not explicitly state that it answers arXiv:1101.5274, but its theorem
does so after the direct locally convex identification above.

\section*{References}

\begin{thebibliography}{9}

\bibitem{BKL}
C. S. Barroso, O. F. K. Kalenda, and P.-K. Lin,
\emph{On the approximate fixed point property in abstract spaces},
arXiv:1101.5274.

\bibitem{MC}
G. Martinez-Cervantes,
\emph{Banach spaces with weak*-sequential dual ball},
arXiv:1612.05948.

\end{thebibliography}

\end{document}
